\chapter*{General Introduction}
\addcontentsline{toc}{chapter}{General Introduction}

\section*{Context of the Global Cybersecurity Landscape}
In the modern digital era, the hyper-connectivity of corporate networks, cloud platforms, and critical infrastructures has created an unprecedented and highly volatile attack surface. Modern organizations rely entirely on software and network availability to operate, which has made them highly attractive targets for malicious actors. The threat landscape has evolved from simple, opportunistic scripts and virus files into highly coordinated, structured, and persistent adversary campaigns~\cite{networks_cybersecurity}. 

These campaigns are led by a diverse array of threat actors, including sophisticated groups conducting digital espionage, ransomware cartels seeking financial extortion, and hacktivists launching disruptive operations. Adversaries routinely exploit security flaws, set up command servers to control attacks, and leak stolen corporate documents. In this hostile landscape, traditional defenses (such as firewalls and antivirus software) are no longer enough to protect sensitive assets.

\section*{The Critical Problematic: Blindness to the Adversary}
A foundational axiom of strategic warfare, articulated by the legendary Chinese military strategist Sun Tzu in \emph{The Art of War}~\cite{art_of_war}, states:
\begin{quote}
    \emph{"If you know the enemy and know yourself, you need not fear the result of a hundred battles. If you know yourself but not the enemy, for every victory gained you will also suffer a defeat. If you know neither the enemy nor yourself, you will succumb in every battle."}
\end{quote}

When applied to modern cybersecurity, the vast majority of security programs focus exclusively on the second clause of Sun Tzu's axiom: \textbf{"knowing themselves"}. Organizations allocate significant resources to internal defense mechanisms. They install detection agents on laptops, aggregate internal logs in monitoring systems, and configure access controls. 

While these tools are essential for monitoring internal systems, they leave the organization blind to \textbf{"the enemy"}. Defenders have no external visibility. They do not know who is targeting their specific sector, which malicious servers are active, what scans are targeting their networks, or if their data is being traded online. Operating in this state leaves security teams entirely reactive, forced to wait until a breach occurs before responding.

\section*{The Solution: Proactive Threat Intelligence (CTI)}
To resolve this imbalance, the field of \textbf{Cyber Threat Intelligence (CTI)} has emerged. CTI is the collection and analysis of data regarding cyber threats. It bridges the visibility gap by providing defenders with \textbf{external awareness}, answering critical questions:
\begin{itemize}
    \item \textbf{Who is the threat?} (Identifying actor profiles and historical campaign records).
    \item \textbf{What are their tactics and methods?} (Understanding adversary techniques using frameworks like the MITRE ATT\&CK matrix, which outlines the phases of an attack).
    \item \textbf{What infrastructure are they using?} (Tracking active malicious servers, phishing domains, and malware download URLs).
\end{itemize}

By integrating threat intelligence, defenders transition from a reactive posture to a proactive defense, identifying and blocking incoming threats before they execute.

\section*{The Metric of Failure: Compromise Dwell Time}
The severe operational cost of defending a network without external threat visibility is illustrated by global security metrics. According to industry reports, the global average **dwell time** -- the duration between an initial compromise and its detection -- consistently exceeds **200 days** (approximately 6 to 7 months). 

During this long window, attackers move through internal networks, gain access to administrative controls, steal intellectual property, and deploy ransomware. The main cause of this long dwell time is that defenders only detect compromises after an internal alert is triggered. By aggregating and streaming external threat intelligence in real-time, organizations can identify attacks immediately, reducing detection time from months to seconds.

\section*{Objective of the Graduation Project}
The primary objective of this graduation project (Projet de Fin d'Études) is to design and implement the **Predictra Threat Map**, a threat intelligence platform. The system aggregates live threat indicators from 9+ external feeds (covering malware sites, malicious servers, and ransomware leaks), finds their geographical locations, classifies target business sectors, and streams these threats in real-time.

On the frontend, these threat flows are displayed on an interactive 3D virtual Earth, allowing teams to instantly see attack vectors. Furthermore, the platform integrates a STIX 2.1~\cite{stix_specification} file parser and a MITRE ATT\&CK mapping grid to support threat analysis.

\section*{Structure of the Graduation Report}
To detail the research, modeling, and realization of the Predictra Threat Map, this report is structured into four distinct chapters:
\begin{itemize}
    \item \textbf{Chapter 1: General Project Framework and Methodology} introduces our host company Predictra, reviews existing threat intelligence platforms, and details our Scrum and UML development methodologies.
    \item \textbf{Chapter 2: Project Specification, Planning, and Architecture} documents functional and non-functional requirements, defines system actors, breaks down the backlog, and maps out the system architecture.
    \item \textbf{Chapter 3: First Release - Core Ingestion and 3D Visualization} details the design and implementation of Sprint 1 (our backend scrapers and target classification engine) and Sprint 2 (our real-time stream and 3D globe visualizer).
    \item \textbf{Chapter 4: Second Release - STIX Analytics and Performance Guardrails} presents Sprint 3 (our STIX 2.1 parser and MITRE matrix) and Sprint 4 (our searchable logs browser, "Target My IP" locator, circular memory buffers, and frame rate optimization engine).
\end{itemize}
